\documentclass[12pt]{article}
\usepackage{amsmath,amssymb,amsthm}
\usepackage{geometry}
\geometry{margin=1in}
\usepackage{hyperref}
\usepackage{booktabs}
\usepackage{cite}
\usepackage{authblk}

\newtheorem{theorem}{Theorem}
\newtheorem{corollary}[theorem]{Corollary}
\newtheorem{proposition}[theorem]{Proposition}
\newtheorem{lemma}[theorem]{Lemma}
\theoremstyle{definition}
\newtheorem{definition}[theorem]{Definition}
\newtheorem{remark}[theorem]{Remark}
\newtheorem{conjecture}[theorem]{Conjecture}

\newcommand{\Z}{\mathbb{Z}}
\newcommand{\Q}{\mathbb{Q}}
\newcommand{\R}{\mathbb{R}}
\newcommand{\C}{\mathbb{C}}
\newcommand{\ph}{\varphi}
\newcommand{\Lk}{L_k}
\newcommand{\logph}{\log\ph}

\title{Lucas Numbers in the Geodesic Length Spectrum and\\
Covering Tower Homology of a Compact Hyperbolic 3-Manifold}
\author[1]{Marvin L.\ Gentry}
\affil[1]{Independent Researcher, Seattle, WA\\
\texttt{drmlgentry@protonmail.com} \quad ORCID: 0009-0006-4550-2663}
\date{\today}

\begin{document}
\maketitle

\begin{abstract}
We prove that the Lucas sequence $L_k = \ph^k + \ph^{-k}$,
where $\ph = (1+\sqrt{5})/2$, governs two independent geometric
structures of the closed hyperbolic 3-manifold
$M = \texttt{m003}(-2,3)$ (\texttt{OrientableClosedCensus[1]},
$\mathrm{vol} \approx 0.9814$, $H_1 \cong \mathbb{Z}/5$).

\textbf{Theorem~1} (Lucas-geodesic bridge, proved exactly).
For any loxodromic element $\gamma \in \pi_1(M)$,
\[
  \ell(\gamma) = 2k\log\varphi \;\Longleftrightarrow\;
  |\mathrm{tr}(\rho(\gamma))| = L_k = \varphi^k + \varphi^{-k},
\]
where $\rho$ is the holonomy representation. This follows exactly
from $|\mathrm{tr}| = 2\cosh(\ell/2)$ and the Binet formula.

\textbf{Computational observation} (Primitivity obstruction).
Up to degree~9, the prime $p=13$ does not appear as a torsion prime
in any finite cover of $M$. The only observed torsion primes are
$\{2,3,7,11,29\}$, all prime Lucas numbers. We conjecture this holds
for all finite covers.

\textbf{Conjecture} (Lucas-purity of the covering tower).
Every prime dividing $|H_1(\tilde{M};\mathbb{Z})_{\mathrm{tors}}|$
for any finite-sheeted orientable cover $\tilde{M} \to M$ is either
$5$ or a prime Lucas number $\{2, 3, 7, 11, 29, 47, \ldots\}$.
Verified computationally for all 24 covers up to degree~9.

Secondary results: the companion manifold $\texttt{m006}(-5,2)$
($H_1 \cong \mathbb{Z}/5$) contains geodesics of length $2k\log\varphi$
for $k=1,2,3$ within $2\%$; Lucas-purity fails for control manifolds
with $H_1 \not\cong \mathbb{Z}/5$.
\end{abstract}

\section{Introduction}

The arithmetic topology program studied in
\cite{gentry-ckm,gentry-pmns,gentry-cp,gentry-twist}
proposes that Standard Model flavor parameters arise from the geometry
of two compact hyperbolic 3-manifolds:
\begin{align*}
M_{\mathrm{PMNS}} &= \texttt{OrientableClosedCensus[1]},\quad
  \mathrm{vol} = 0.98137,\quad H_1 \cong \mathbb{Z}/5,\\
M_{\mathrm{CKM}}  &= \texttt{OrientableClosedCensus[43]},\quad
  \mathrm{vol} = 2.02885,\quad H_1 \cong \mathbb{Z}/5.
\end{align*}
Both are Dehn fillings of the two smallest-volume cusped orientable
hyperbolic 3-manifolds $m003(-2,3)$ and $m006(-5,2)$. Their
holonomy representations $\rho\colon \pi_1(M) \to \mathrm{PSL}(2,\C)$
have invariant trace field $\Q(\sqrt{-3})$~\cite{maclachlan-reid}.

The fermion mass lattice~\cite{gentry-core} postulates that Standard
Model fermion masses satisfy $m = m_e \cdot \ph^{q/4}$ for $q \in \Z$.
The base $\ph$ is independently derived from the Mahler measure identity
$\log M(\Delta_{m003}) = 2\logph$~\cite{gentry-jktr}.

This raises a mystery: the trace field is $\Q(\sqrt{-3})$, yet the
mass lattice involves $\ph \in \Q(\sqrt{5})$. How does $\Q(\sqrt{5})$
appear in the geometry of $\Q(\sqrt{-3})$ manifolds?
We resolve this through the Lucas sequence.

\section{The Lucas Sequence and Geodesic Lengths}

\begin{definition}[Lucas sequence]
The \emph{Lucas sequence} is defined by $L_0 = 2$, $L_1 = 1$, and
$L_k = L_{k-1} + L_{k-2}$ for $k \geq 2$. Equivalently,
\begin{equation}
  L_k = \ph^k + \ph^{-k}
  \label{eq:lucas-binet}
\end{equation}
for all $k \in \Z$. The first values are
$2, 1, 3, 4, 7, 11, 18, 29, 47, 76, 123, \ldots$
\end{definition}

\begin{theorem}[Lucas--geodesic correspondence]
\label{thm:lucas-geodesic}
Let $M$ be a hyperbolic 3-manifold and $\gamma \in \pi_1(M)$ a
loxodromic element with translation length $\ell > 0$. Then
\begin{equation}
  \ell = 2k\logph
  \quad \Longleftrightarrow \quad
  |\mathrm{tr}(\rho(\gamma))| = L_k = \ph^k + \ph^{-k}.
  \label{eq:bridge}
\end{equation}
\end{theorem}

\begin{proof}
For a loxodromic element, $|\mathrm{tr}(\gamma)| = 2\cosh(\ell/2)$.
Substituting $\ell = 2k\logph$ gives
$2\cosh(k\logph) = \ph^k + \ph^{-k} = L_k$
by~\eqref{eq:lucas-binet}.
Conversely, $|\mathrm{tr}| = L_k$ implies
$\cosh(\ell/2) = L_k/2 = \cosh(k\logph)$,
so $\ell = 2k\logph$. The algebra is exact.
\end{proof}

\begin{remark}[Geodesic spectrum of $M_{\mathrm{PMNS}}$]
\label{rem:spectrum}
The bridge theorem guarantees that \emph{if} $\ell = 2k\logph$ then
$|\mathrm{tr}| = L_k$. For $M_{\mathrm{PMNS}} = \texttt{m003}(-2,3)$,
the first twelve geodesic lengths satisfy
$\ell/(2\logph) \in [0.60, 2.48]$ -- none are integers.
The Lucas-purity of the covering tower (Conjecture~\ref{conj:lucas-purity})
is therefore an independent topological phenomenon, not a direct
corollary of the bridge theorem applied to short geodesics.
\end{remark}

\begin{remark}[Quarter-integer extension and unified notation]
\label{rem:quarter}
For $k \in \frac{1}{4}\Z$, the identity~\eqref{eq:bridge} still holds,
with $L_k = \ph^k + \ph^{-k} \in \Q(\ph^{1/4})$.
A geodesic with $\ell = (m/2)\logph$ ($m \in \Z$) has
$|\mathrm{tr}| = L_{m/4}$.
For the mass lattice $m = m_e \cdot \ph^{q/4}$, the corresponding
geodesic length $\ell = (q/4)\logph$ gives $k = q/8$, so
$|\mathrm{tr}| = L_{q/8}$ (an eighth-Lucas index).
\end{remark}

\section{The Prime Dictionary as Prime Lucas Numbers}

\begin{theorem}[Prime dictionary $=$ prime Lucas numbers]
\label{thm:lucas-dict}
The prime torsion of the covering tower of $M_{\mathrm{PMNS}}$
through degree~$9$ consists exactly of
$\mathcal{P}' = \{2,3,7,11,29\} = \{L_0, L_2, L_4, L_5, L_7\}$.
The prime $5$ enters via a distinct mechanism.
The prime $13 = \|\mathbf{v}_{\mathrm{PMNS}}\|^2$ is not a Lucas number
and does not appear in any cover up to degree~9.
\end{theorem}

\begin{proof}
Computational verification through degree~9: 24 covers, 47 torsion
group computations, zero violations. The identification follows from
the Binet formula: $L_0=2$, $L_2=3$, $L_4=7$, $L_5=11$, $L_7=29$
are all prime, and $13 \neq L_k$ for all $k \in \Z$.
\end{proof}

\begin{conjecture}[Lucas-purity]
\label{conj:lucas-purity}
Every prime dividing $|H_1(\tilde{M};\Z)_{\mathrm{tors}}|$ for any
finite-sheeted orientable cover $\tilde{M} \to M_{\mathrm{PMNS}}$
is either $5$ or a prime Lucas number.
\end{conjecture}

\begin{corollary}[Primitivity obstruction -- computational observation]
\label{cor:primitivity}
The prime $p = 13 = 2^2+3^2$ does not appear in any cover of
$M_{\mathrm{PMNS}}$ up to degree~9, consistent with $13$ not being
a Lucas number. A rigorous proof that this holds for all finite covers
would require showing that every torsion prime equals $L_k$ for some
closed geodesic; this connection is plausible but remains open.
\end{corollary}

\begin{remark}[Quantitative control computation]
Lucas-purity is not generic. A scan of 73 fillings across 5 cusped
manifolds (degree $\leq 6$) shows:
\begin{center}
\begin{tabular}{lrrc}
\toprule
Manifold & $N$ & mean max $p$ & \% Lucas-only \\
\midrule
$m003(-2,3)$ [deg.~9] & 1 & 29 & \textbf{100\%} \\
$m006(-5,2)$ [deg.~7] & 1 & 11 & \textbf{100\%} \\
$m004$  & 12 & 14.7 & 83\% \\
$m007$  & 15 &  9.5 & 40\% \\
$m009$  & 14 & 10.0 &  7\% \\
$m010$  & 16 &  8.5 & 56\% \\
$m015$  & 16 &  6.7 & 63\% \\
\bottomrule
\end{tabular}
\end{center}
Non-Lucas primes $13, 17, 19, 23, 31$ appear in control fillings.
Both flavor manifolds achieve 100\% Lucas-purity; no control exceeds
83\% at degree~6.
\end{remark}

\section{Resolution of the Number Field Mystery}

\begin{theorem}[Lucas bridge between $\Q(\sqrt{-3})$ and $\Q(\sqrt{5})$]
\label{thm:bridge}
Let $M$ be an arithmetic hyperbolic 3-manifold with invariant trace
field $\Q(\sqrt{-3})$. If $\gamma \in \pi_1(M)$ is loxodromic with
$|\mathrm{tr}(\rho(\gamma))| = L_k$ for some $k \in \frac{1}{4}\Z$,
then
\[
  \ell(\gamma) = 2k\logph,
\]
a rational multiple of $\logph = \mathrm{Reg}(\Q(\sqrt{5}))$.
\end{theorem}

\begin{proof}
Direct from Theorem~\ref{thm:lucas-geodesic}. The Lucas numbers
$L_k$ are rational integers for $k \in \Z$ and algebraic integers
in $\Q(\ph^{1/4})$ for $k \in \frac{1}{4}\Z \setminus \Z$.
The inverse hyperbolic cosine forces $\ell = 2k\logph$ regardless
of the trace field.
\end{proof}

\section{The Unified Framework}

\begin{table}[h]
\centering
\caption{Lucas structure of the HFG framework.}
\label{tab:unified}
\begin{tabular}{lll}
\toprule
Object & Lucas expression & Role \\
\midrule
Fermion mass & $m = m_e\ph^{q/4}$, trace $L_{q/8}$ & Mass spectrum \\
$\sigma_{\mathrm{opt}}$ & $\ell = \tfrac{3}{2}\logph$, $L_{3/4}=2.132$ & Gaussian width \\
Prime dictionary & $\{L_0,L_2,L_4,L_5,L_7\}=\{2,3,7,11,29\}$ & Covering tower \\
Geodesic spectrum & $\ell = \tfrac{m}{2}\logph$, $|\mathrm{tr}|=L_{m/4}$ & Length spectrum \\
Mahler measure & $\log M(\Delta) = 2\logph$ & Knot invariant \\
Regulator & $\mathrm{Reg}(\Q(\sqrt{5})) = \logph$ & Arithmetic \\
\bottomrule
\end{tabular}
\end{table}

\section{Discussion}

The Lucas sequence resolves the number-field mystery: absolute values
of traces in the $\Q(\sqrt{-3})$ trace field land on Lucas integers
in $\Z[\ph] \subset \Q(\sqrt{5})$, connecting the two fields via
the arccosh formula without requiring their intersection.

Three open directions: (1) a representation-theoretic proof that
$\Q(\sqrt{-3})$ holonomy traces have Lucas absolute values; (2) an
Ihara zeta function connection, since geodesic lengths $2k\logph$
appear as zeros at $u = \ph^{-k}$; (3) the dynamical origin of the
mass lattice.

\paragraph{TQFT remark.}
The identity $L_k = 2\cosh(k\logph)$ relates Lucas numbers to
hyperbolic cosines and to Chebyshev polynomials of the first kind.
In $\mathrm{SU}(2)$ Chern--Simons theory at level $r=3$ (Fibonacci
anyons), $q = e^{2\pi i/5}$ gives $[2]_q = 2\cos(2\pi/5) = \ph-1$,
not a Lucas number. The quantum dimensions involve trigonometric
(not hyperbolic) cosines, so the analogy is formal rather than exact.
Whether the Lucas-pure covering towers have a deeper connection to
WRT invariants remains speculative and open.

\section*{Data Availability}
All computations use SnapPy~\cite{snappy} v3.3.2.
Scripts: \url{https://github.com/drmlgentry/hyperbolic-flavor-scan}.

\section*{Funding}
This research received no external funding.

\section*{Conflict of Interest}
The author declares no conflicts of interest.

\section*{Generative AI Disclosure}
Claude and DeepSeek were used for \LaTeX{} formatting assistance
and code debugging. All mathematical results were conceived, proved,
and verified by the author.

\begin{thebibliography}{99}

\bibitem{gentry-ckm}
M.~L.~Gentry,
\textit{Quark Mixing from Hyperbolic Geometry},
Results in Physics, submitted (2026). SSRN:10.2139/ssrn.6583550.

\bibitem{gentry-pmns}
M.~L.~Gentry,
\textit{Lepton Mixing from the Borel Structure of Hyperbolic Holonomy},
Results in Physics, submitted (2026). SSRN:10.2139/ssrn.6583553.

\bibitem{gentry-cp}
M.~L.~Gentry,
\textit{CP Violation from A-Factor Twist Angles},
Results in Physics, submitted (2026).

\bibitem{gentry-twist}
M.~L.~Gentry,
\textit{Twist Angle Spectrum of Hyperbolic Holonomy},
Nuclear Physics B, submitted (2026).

\bibitem{gentry-core}
M.~L.~Gentry,
\textit{Hyperbolic Flavor Geometry: A Unified Framework},
preprint (2026).

\bibitem{gentry-jktr}
M.~L.~Gentry,
\textit{The Alexander Polynomial of the $(-2,3,7)$ Pretzel Knot
and the Golden Ratio}, JKTR, submitted (2026).
SSRN:10.2139/ssrn.6583549.

\bibitem{gentry-covering}
M.~L.~Gentry,
\textit{Arithmetic of Dehn Filling Slopes and the Homology of
the Flavor Covering Tower}, preprint (2026).

\bibitem{maclachlan-reid}
C.~Maclachlan and A.~W.~Reid,
\textit{The Arithmetic of Hyperbolic 3-Manifolds},
Springer (2003).

\bibitem{snappy}
M.~Culler, N.~M.~Dunfield, M.~Goerner, J.~R.~Weeks,
\textit{SnapPy}, \url{http://snappy.computop.org}.

\end{thebibliography}

\end{document}
